% Chapter 8: risks with likelihood, impact and mitigation.
\chapter{Risks and mitigation}\label{ch:risks}

Table~\ref{tab:risks} lists the risks judged most likely to delay the
thesis. Likelihood and impact are rated low, medium or high; the
mitigation column states what is already in place.

\newcommand*{\rate}[1]{\textcolor{accent}{#1}}
\begin{table}[ht]
\centering\small
\caption{Main risks to the programme.}
\label{tab:risks}
\begin{tabularx}{\linewidth}{@{}>{\raggedright\arraybackslash}p{36mm}
  >{\sffamily\footnotesize}c >{\sffamily\footnotesize}c
  >{\raggedright\arraybackslash}X@{}}
\toprule
\thead{Risk} & \thead{Likel.} & \thead{Impact} & \thead{Mitigation}\\
\midrule
Mounting permits are delayed beyond May 2027
  & \rate{med} & high
  & List of 140 pre-approved lampposts agreed with the city in 2026;
    campaign can start with these alone. \\
A cool summer in 2027 gives few heat days
  & \rate{med} & med
  & Nights above 20\,°C are frequent even in cool summers; the
    analysis also uses the 2028 municipal data. \\
Partner city withdraws data access
  & low & high
  & Data-sharing agreements signed with both cities; Westhaven data are
    already archived. \\
Sensor loss or vandalism above 20\,\%
  & \rate{med} & low
  & Ten spare units; tamper-resistant mounts; loss is itself recorded
    as a result for \ref{rq:field}. \\
Weighted rule shows no gain over variance-only
  & low & med
  & A null result is publishable and still answers \ref{rq:acq};
    Chapter~5 does not depend on it. \\
\bottomrule
\end{tabularx}
\end{table}
